\scanpage{17}
\mspage{7}
Le groupe fondamental mixte est défini à la \unsure{Malgoire}, \deleted{\unsure{comprenant}} les classes $(\gamma,l)$ d'un $\gamma\in\Gamma$ et d'une classe de chemins $l$ de $P$ dans $\gamma^{-1}(P)$, en définissant
\[
(\gamma,l)(\gamma',l')=(\inserted{\gamma\gamma',}\gamma'^{-1}(l)\circ l')
\]
i.e. ``composant les chemins comme des fonctions'' (c'est la convention pour faire opérer à \uline{gauche} le groupoïde fondamental sur les fibres d'un rev. étale \inserted{$U'$}, à une classe de chemins $l$ de $x$ à $y$ étant associé $l:U'_x\longrightarrow U'_y\ldots$).

On choisit des éléments
\[
\sigma_0,\sigma_1,\sigma_\infty\qquad\text{sur}\qquad\dot\sigma_0,\dot\sigma_1,\dot\sigma_\infty
\]
comme les couples d'un $\unsure{\sigma_i|\dot\sigma_i}$, et des chemins \inserted{$\lambda_i$} de $P$ à $\bar P=\unsure{\sigma_i|\dot\sigma_i}P$ qui se fait sur \deleted{l'unique} \inserted{celui des} demi-grands cercles de $P$ à $\bar P$, \deleted{stable par la symétrie} \inserted{\unsure{ens. des pts fixes}} \deleted{$\tau\sigma_i$} de la symétrie $\dot\tau\dot\sigma_i=\dot\sigma_i\dot\tau$ (symétrie \unsure{par un plan}, fixant $P,\bar P,a_i$ $(=0,1\text{ ou }\infty)$ et échangeant les deux autres $a_j,a_k$), et \unsure{ne} passant pas par $a_i$. \deleted{\unsure{Au}} Désignant par $l_0,l_1,l_\infty$ \inserted{(\unsure{sur|pour} $1\in\Gamma$)} les lacets en sens direct autour des $0,1,\infty$ à partir de $P$, en allant vers $0,1,\infty$ suivant des \unsure{quarts} de grands cercles, on trouve
\[
(1)\qquad\big|\quad l_\infty l_1l_0=1
\]
(et non $l_0l_1l_\infty=1$ ! \unsure{?} On peut l'écrire $l_0^{-1}l_1^{-1}l_\infty^{-1}=1$)
et les relations
\mspage{16}
